\chapter{Gaussian Processes and the Birth of Brownian Motion}

\begin{center}\itshape
A random function is a family of random numbers, one for each time. Kolmogorov showed that any consistent way of choosing\\
finitely many of them at once defines the whole family. But continuity is a property of uncountably many times,\\
and to see it one needs one more idea: control the increments, and the path will hold together.
\end{center}

In the summer of 1827 the Scottish botanist Robert Brown put grains of pollen of the wild flower \emph{Clarkia pulchella} in water under his microscope and saw
the tiny particles that came out of them jiggling without rest. He suspected life, so he tried dead material: pollen kept for a century in a herbarium, then
powdered glass, granite, and even a fragment of the Sphinx. Everything small enough danced. He published his ``brief account of microscopical observations'' the next year,
and could not explain what he had seen.

The explanation took eighty years. In 1905 Albert Einstein, then a clerk at the patent office in Bern, argued that if water is made of molecules in thermal agitation,
a suspended particle must be kicked by them at random, and its displacement over a time $t$ must have a mean square proportional to $t$, not to $t^2$ as for
a steady motion. The constant of proportionality involves the number of molecules in a mole, Avogadro's number. So by watching a grain of resin under a
microscope one could count molecules. Jean Perrin did exactly that in 1908 and 1909, following particles of gamboge by hand with a camera lucida, and his
values of Avogadro's number agreed with those found by entirely different methods. The reality of atoms, still doubted by eminent chemists in 1900, was no longer
in question, and Perrin received the Nobel Prize in 1926. In his book \emph{Les Atomes} (1913) Perrin drew the paths of his particles, and remarked that the true
path, looked at more and more closely, would reveal ever finer zigzags: it is the kind of curve without tangents that mathematicians had invented as a curiosity
(this is a paraphrase of his preface). Meanwhile, in 1900, Louis Bachelier had used the same process in his thesis on the speculation in the Paris stock
exchange, with nobody noticing for half a century.

The first mathematical construction came from Norbert Wiener in 1923. Wiener built a measure on the space of continuous functions under which the increments are
independent Gaussians, the measure that probabilists now call \emph{Wiener measure}. His methods were those of the integration theory of the time, and they were
heavy. Ten years later Kolmogorov's \emph{Grundbegriffe} gave the general framework: a random process is a probability on a space of functions, determined by its
finite-dimensional laws. But Kolmogorov's measure lives on all functions, and the set of continuous functions is not even an event for it. The bridge was a
criterion that Kolmogorov found and Evgeny Slutsky published in 1937: if increments are small in mean, at the right rate, the process has a continuous version.
And Wiener's own favourite description, written with Raymond Paley and Antoni Zygmund in 1933, is the most concrete of all: Brownian motion is a Fourier series
with independent Gaussian coefficients,
\[B_t=\xi_0t+\sqrt2\sum_{n\ge1}\xi_n\frac{\sin n\pi t}{n\pi}.\]

Tome~3 built Brownian motion by L\'evy's midpoint method and studied its first properties. This chapter builds it twice more, in the two ways just described,
because each construction reveals something the other hides: Kolmogorov's criterion gives the H\"older regularity of the path, and the Fourier series shows
that Brownian motion is a sum of independent sine waves. Along the way we meet the general Gaussian process and its most famous relative after Brownian motion,
the fractional Brownian motion of Kolmogorov and Mandelbrot, whose increments remember the past.

\begin{reminder}[from Tomes 3, 4, 6 and 7]
\emph{Gaussian vectors} (Tome~3, Theorem~6.15): a random vector $X$ of $\R^d$ is Gaussian if every linear form $\langle u,X\rangle$ is Gaussian. Its law is determined by its mean $m$ and covariance
$\Sigma$, since its characteristic function is $e^{i\langle u,m\rangle-u^{\mathsf T}\Sigma u/2}$; every symmetric positive semi-definite $\Sigma$ is the covariance of a Gaussian vector; uncorrelated components of a
Gaussian vector are independent; and $AX+b$ is Gaussian. \emph{Brownian motion} (Tome~3, Definition~11.1 and Proposition~11.2): a process with $B_0=0$, continuous paths and independent Gaussian increments
$B_t-B_s\sim N(0,t-s)$; equivalently a centred Gaussian process with continuous paths and covariance $\min(s,t)$. It exists (Tome~3, Theorem~11.4), it is nowhere differentiable (Tome~3,
Theorem~11.10), and random walks converge to it (Tome~10, Theorem~3.13). \emph{Kolmogorov's extension theorem} (Tome~6, Theorem~11.8): consistent probabilities $\mu_n$ on $\R^n$ (the image of
$\mu_{n+1}$ by the projection forgetting the last coordinate is $\mu_n$) come from a unique probability on $\R^\N$ with its product $\sigma$-algebra. \emph{Parseval} (Tome~4, Theorem~6.11): for
piecewise continuous $2\pi$-periodic $f,g$, $\frac1{2\pi}\int_{-\pi}^{\pi}f\bar g=\sum_{k\in\Z}c_k(f)\overline{c_k(g)}$. \emph{Density} (Tome~6, Theorem~9.7): continuous functions with compact support are dense in $L^2(\R)$; since such a function is a uniform limit of step functions supported in a fixed compact set,
step functions are dense too. The \emph{Karhunen--Lo\`eve modes} of Brownian motion on $[0,1]$ (Tome~7, Chapter~4): $\sqrt2\sin((n-\frac12)\pi t)$, with eigenvalues $((n-\frac12)\pi)^{-2}$.
\end{reminder}

\section{Gaussian processes}

A \emph{process} indexed by a set $T$ is a family $(X_t)_{t\in T}$ of random variables on one probability space. Its \emph{finite-dimensional laws} are the laws of the vectors $(X_{t_1},\dots,X_{t_k})$. It is
\emph{Gaussian} if all of them are Gaussian, and then everything is determined by the mean $m(t)=\mathbb E X_t$ and the covariance $K(s,t)=\operatorname{Cov}(X_s,X_t)$. Such a $K$ is symmetric and
\emph{positive semi-definite}: $\sum_{j,k}c_jc_kK(t_j,t_k)=\operatorname{Var}\bigl(\sum_jc_jX_{t_j}\bigr)\ge0$ for all finite families. Does every such kernel come from a process? Kolmogorov's theorem
answers yes, once it is extended from sequences to arbitrary index sets.

For $F\subseteq G\subseteq T$ let $\pi^G_F:\R^G\to\R^F$ be the restriction of functions, and $\pi_F=\pi^T_F$. On $\R^T$ we use the \emph{product $\sigma$-algebra}, generated by the coordinates
$x\mapsto x(t)$.

\begin{theorem}[Kolmogorov's extension theorem, any index set]\label{thm:kolmext}
Let $T$ be a set and, for every finite $F\subseteq T$, let $\mu_F$ be a probability on $\R^F$. Suppose the family is \emph{consistent}: for $F\subseteq G$ finite, the image of $\mu_G$ by $\pi^G_F$ is $\mu_F$.
Then there is a unique probability $\mu$ on $\R^T$ whose image by $\pi_F$ is $\mu_F$ for every finite $F$.
\end{theorem}

\begin{proof}
\emph{Countable index sets.} If $S=\{s_1,s_2,\dots\}\subseteq T$ is countable, the laws $\mu_{\{s_1,\dots,s_n\}}$, read on $\R^n$, are consistent in the sense of Tome~6, Theorem~11.8, which gives a probability $\mu_S$ on
$\R^S\cong\R^\N$ with the right images on the sets $\{s_1,\dots,s_n\}$, and then on every finite $F\subseteq S$ (contained in some $\{s_1,\dots,s_n\}$). By uniqueness (Tome~6, Theorem~6.5, the cylinder sets
forming a $\pi$-system), $\mu_S$ does not depend on the enumeration, and for countable $S\subseteq S'$ the image of $\mu_{S'}$ by $\pi^{S'}_S$ is $\mu_S$, since both have the images $\mu_F$.

\emph{Events depend on countably many coordinates.} Let $\mathcal A$ be the family of sets $\pi_S^{-1}(A)$ with $S\subseteq T$ countable and $A$ in the product $\sigma$-algebra of $\R^S$. It contains the
coordinate events, it is stable under complement, and under countable unions: given $\pi_{S_n}^{-1}(A_n)$, let $S=\bigcup S_n$, still countable, and note $\pi_{S_n}^{-1}(A_n)=\pi_S^{-1}\bigl((\pi^S_{S_n})^{-1}A_n\bigr)$.
So $\mathcal A$ is the whole product $\sigma$-algebra.

\emph{Definition of $\mu$.} Put $\mu(\pi_S^{-1}A)=\mu_S(A)$. If $\pi_S^{-1}A=\pi_{S'}^{-1}A'$, let $S''=S\cup S'$; since $\pi_{S''}$ is onto, $(\pi^{S''}_S)^{-1}A=(\pi^{S''}_{S'})^{-1}A'$, and both sides have
$\mu_{S''}$-measure $\mu_S(A)=\mu_{S'}(A')$ by the consistency just proved. So $\mu$ is well defined. Countably many events $\pi_{S_n}^{-1}(A_n)$ can be written over the single countable set
$S=\bigcup S_n$, where $\mu$ is the measure $\mu_S$ transported; so $\mu$ is countably additive. Uniqueness: two such measures agree on the $\pi$-system of finite-dimensional cylinders, which
generates the product $\sigma$-algebra.
\end{proof}

\begin{corollary}[Gaussian processes exist]\label{cor:gaussexist}
For every set $T$, every function $m:T\to\R$ and every symmetric positive semi-definite kernel $K$ on $T$, there is a Gaussian process indexed by $T$ with mean $m$ and covariance $K$, and its
law on $\R^T$ is unique.
\end{corollary}

\begin{proof}
For finite $F$, let $\mu_F$ be the Gaussian law on $\R^F$ with mean $(m(t))_{t\in F}$ and covariance matrix $(K(s,t))_{s,t\in F}$, which is positive semi-definite (Tome~3, Theorem~6.15(b)). The restriction
$\pi^G_F$ is linear, so the image of $\mu_G$ is Gaussian with the restricted mean and the sub-matrix of covariances (Theorem~6.15(d)), that is $\mu_F$. Apply Theorem~\ref{thm:kolmext} and take $X_t(x)=x(t)$ on
$\R^T$. Uniqueness: two such laws agree on cylinders.
\end{proof}

\begin{example}\label{ex:kernels}
Four kernels on $[0,\infty[$, each positive semi-definite because it is the covariance of something we can build by hand.
\begin{enumerate}[label=(\alph*)]
\item $\min(s,t)$: $\min(s,t)=\int_0^\infty\mathbf 1_{[0,s]}\mathbf 1_{[0,t]}$, so $\sum c_jc_k\min(t_j,t_k)=\int\bigl(\sum c_j\mathbf 1_{[0,t_j]}\bigr)^2\ge0$. This is \emph{Brownian motion}, as far as the finite-dimensional laws go.
\item $\min(s,t)-st$ on $[0,1]$: the covariance of $B_t-tB_1$, the \emph{Brownian bridge}.
\item $e^{-|t-s|}$: the covariance of the \emph{stationary Ornstein--Uhlenbeck process} $e^{-t}W_{e^{2t}}$ for a Brownian motion $W$, as $\operatorname{Cov}=e^{-s-t}\min(e^{2s},e^{2t})=e^{-|t-s|}$.
\item $\frac12(s^{2H}+t^{2H}-|t-s|^{2H})$ with $0<H<1$: \emph{fractional Brownian motion}, built in Section~\ref{sec:fbm}.
\end{enumerate}
\end{example}

The process of Corollary~\ref{cor:gaussexist} lives on $\R^T$, and here lies a real difficulty when $T=[0,1]$.

\begin{think}
Is the set $C[0,1]$ of continuous functions an event of the product $\sigma$-algebra of $\R^{[0,1]}$? (By the proof of Theorem~\ref{thm:kolmext}, every event is determined by the values at countably
many points.)
\end{think}

It is not: an event $\pi_S^{-1}(A)$ with $S$ countable contains, together with any continuous function, all functions that agree with it on $S$, among them discontinuous ones. So ``the path is
continuous'' has no probability under Kolmogorov's measure, and in fact nothing can be said about it from the finite-dimensional laws alone: if $U$ is uniform on $[0,1]$, the processes
$X_t=0$ and $Y_t=\mathbf 1_{\{t=U\}}$ have the same finite-dimensional laws, and one is continuous while the other never is. The way out is to change each $X_t$ on a null set.

\begin{definition}\label{def:modification}
A process $\tilde X$ is a \emph{modification} of $X$ (on the same space) if $P(\tilde X_t=X_t)=1$ for every $t$. Then $X$ and $\tilde X$ have the same finite-dimensional laws.
\end{definition}

\section{Kolmogorov's continuity criterion}

The idea is to look at the process only at dyadic points first, where there are countably many, and to show that it is uniformly continuous there.

\begin{theorem}[Kolmogorov and Chentsov]\label{thm:kolmcont}
Let $(X_t)_{t\in[0,1]^d}$ be a real process, and suppose there are $p,\beta,C>0$ with
\[
\mathbb E|X_t-X_s|^p\le C|t-s|^{d+\beta}\qquad\text{for all }s,t\in[0,1]^d,
\]
where $|\cdot|$ is the max norm. Then $X$ has a modification $\tilde X$ whose paths are H\"older continuous of every order $\gamma<\beta/p$: for each such $\gamma$ there is a random $K_\gamma<\infty$ with
$|\tilde X_t-\tilde X_s|\le K_\gamma|t-s|^\gamma$ for all $s,t$.
\end{theorem}

\begin{proof}
Fix $\gamma<\beta/p$. Let $D_n=2^{-n}\Z^d\cap[0,1]^d$, the dyadic points of level $n$, and $D=\bigcup_nD_n$. Call $s,t\in D_n$ \emph{neighbours} if $|t-s|\le2^{-n}$. There are at most $(2^n+1)^d3^d$ ordered pairs
of neighbours, so by Markov's inequality
\[
P\Bigl(\max_{\text{neighbours }s,t\in D_n}|X_t-X_s|>2^{-\gamma n}\Bigr)\le(2^n+1)^d3^d\,C2^{-n(d+\beta)}2^{\gamma pn}\le C6^d2^{-n(\beta-\gamma p)}.
\]
The right side is summable, so by Borel--Cantelli there is an almost surely finite $N$ with $|X_t-X_s|\le2^{-\gamma n}$ for all neighbours $s,t\in D_n$ and all $n\ge N$.

\emph{Chaining.} Let $m\ge N$ and $s,t\in D$ with $|t-s|\le2^{-m}$. For $n\ge m$ let $s_n=2^{-n}\lfloor2^ns\rfloor$ (coordinatewise), and similarly $t_n$. Each coordinate of $2^ns_n$ and $2^nt_n$ differs by at most
$1$, because the coordinates of $2^ns$ and $2^nt$ differ by at most $1$; so $s_m,t_m$ are neighbours in $D_m$. And $s_n,s_{n+1}$ are neighbours in $D_{n+1}$, each coordinate of $s_{n+1}-s_n$ being $0$ or
$2^{-n-1}$. Since $s,t\in D$, $s_n=s$ and $t_n=t$ for $n$ large. Telescoping,
\[
|X_t-X_s|\le|X_{t_m}-X_{s_m}|+\sum_{n\ge m}\bigl(|X_{s_{n+1}}-X_{s_n}|+|X_{t_{n+1}}-X_{t_n}|\bigr)\le2^{-\gamma m}+2\sum_{n>m}2^{-\gamma n}\le c_\gamma2^{-\gamma m},
\]
with $c_\gamma=1+\frac2{2^\gamma-1}$. Now let $s,t\in D$ with $0<|t-s|\le2^{-N}$, and choose $m\ge N$ with $2^{-m-1}<|t-s|\le2^{-m}$: then $|X_t-X_s|\le c_\gamma2^{-\gamma m}\le2^\gamma c_\gamma|t-s|^\gamma$.
Being uniformly continuous on the dense subset $D$ of the compact cube, $X|_D$ is bounded, by $R$ say, and for $|t-s|>2^{-N}$, $|X_t-X_s|\le2R\le2R\,2^{\gamma N}|t-s|^\gamma$. So on the event
$\{N<\infty\}$ the restriction of $X$ to $D$ is H\"older of order $\gamma$, with the random constant $K_\gamma=\max(2^\gamma c_\gamma,2R2^{\gamma N})$.

\emph{Extension.} On $\{N<\infty\}$, $X|_D$ is uniformly continuous on the dense set $D$, so it has a unique continuous extension $\tilde X$ to $[0,1]^d$, still $\gamma$-H\"older with the same constant; on the null
set $\{N=\infty\}$ put $\tilde X=0$. Each $\tilde X_t$ is a random variable, as a pointwise limit of $X_{t_n}$ with $t_n\in D$ (and $0$ on an event). For $t\in D$, $\tilde X_t=X_t$ almost surely. For other $t$, take $t_n\in D$ with
$t_n\to t$: $X_{t_n}\to\tilde X_t$ almost surely, and $X_{t_n}\to X_t$ in probability since $P(|X_{t_n}-X_t|>\varepsilon)\le C\varepsilon^{-p}|t_n-t|^{d+\beta}$. A sequence converging in probability has a subsequence
converging almost surely, so $\tilde X_t=X_t$ almost surely. Finally the event $\{N<\infty\}$ and the constant $K_\gamma$ can be chosen to work for $\gamma=\gamma_j$ along a sequence $\gamma_j\uparrow\beta/p$, hence for all
$\gamma<\beta/p$ at once.
\end{proof}

The proof deserves a second look, because it is the prototype of \emph{chaining}: a supremum over a continuum is bounded by a sum, over scales, of maxima over finite sets, and at each scale a union bound
suffices because the increments get small fast enough to beat the number of pairs.

\begin{corollary}[Brownian motion, a second time]\label{cor:bmholder}
There is a Brownian motion $(B_t)_{t\ge0}$, and almost surely its paths are H\"older continuous of every order $\gamma<\frac12$ on every bounded interval.
\end{corollary}

\begin{proof}
Let $X$ be the centred Gaussian process with covariance $\min(s,t)$ on $[0,\infty[$ (Example~\ref{ex:kernels}(a)). For $s<t$, $X_t-X_s$ is centred Gaussian with variance $t-2s+s=t-s$, so
$\mathbb E|X_t-X_s|^{2k}=(2k-1)!!\,|t-s|^k$ for every integer $k\ge1$. On $[j,j+1]$ (a translate of $[0,1]$), Theorem~\ref{thm:kolmcont} with $d=1$, $p=2k$, $\beta=k-1$ gives a continuous modification, H\"older of
every order $<\frac{k-1}{2k}$; letting $k\to\infty$, of every order $<\frac12$. Two continuous modifications on overlapping intervals agree almost surely at the rational points, hence everywhere, so they glue into a
continuous modification on $[0,\infty[$, with $B_0=X_0=0$ almost surely (we may force $B_0=0$ by setting $B=0$ on a null event). A centred Gaussian process with continuous paths and covariance $\min(s,t)$ is a
Brownian motion (Tome~3, Proposition~11.2).
\end{proof}

Order $\frac12$ itself is out of reach: Tome~10, Exercise~3.7 shows that $\sup_{|t-s|\le h}|B_t-B_s|/\sqrt h$ is unbounded, and Chapter~2 finds the exact modulus of continuity $\sqrt{2h\log(1/h)}$.
The same theorem in dimension two builds the \emph{Brownian sheet}, a random surface (Exercise~\ref{ex:sheet}).

\section{Brownian motion as a Fourier series}

Here is a third construction, in the spirit of Paley, Wiener and Zygmund. Let $(e_n)$ be an orthonormal family in $L^2[0,1]$ and $(\xi_n)$ independent standard Gaussians. The candidate is
\[
W_t=\sum_n\xi_n\int_0^te_n(u)\dd u=\sum_n\xi_n\langle\mathbf 1_{[0,t]},e_n\rangle,
\]
a sum of independent Gaussians with variances $\langle\mathbf 1_{[0,t]},e_n\rangle^2$, whose sum is at most $\|\mathbf 1_{[0,t]}\|^2=t$ by Bessel's inequality. So the series converges in $L^2$ for each $t$ (the partial sums
form a Cauchy sequence, their increments being orthogonal). If the family is complete,
\[
\operatorname{Cov}(W_s,W_t)=\sum_n\langle\mathbf 1_{[0,s]},e_n\rangle\langle\mathbf 1_{[0,t]},e_n\rangle=\langle\mathbf 1_{[0,s]},\mathbf 1_{[0,t]}\rangle=\min(s,t):
\]
the right covariance. We need two concrete orthonormal families, and only the identity of covariances for indicator functions, which Parseval's theorem provides.

\begin{lemma}\label{lem:twobases}
For $s,t\in[0,1]$:
\begin{enumerate}[label=(\alph*)]
\item with $c_0=1$ and $c_n(u)=\sqrt2\cos n\pi u$ ($n\ge1$):
\[\sum_{n\ge0}\langle\mathbf 1_{[0,s]},c_n\rangle\langle\mathbf 1_{[0,t]},c_n\rangle=st+2\sum_{n\ge1}\frac{\sin n\pi s\,\sin n\pi t}{n^2\pi^2}=\min(s,t);\]
\item with $k_n(u)=\sqrt2\cos\bigl((n-\frac12)\pi u\bigr)$ ($n\ge1$):
\[\sum_{n\ge1}\langle\mathbf 1_{[0,s]},k_n\rangle\langle\mathbf 1_{[0,t]},k_n\rangle=\sum_{n\ge1}\frac{2\sin((n-\frac12)\pi s)\sin((n-\frac12)\pi t)}{(n-\frac12)^2\pi^2}=\min(s,t).\]
\end{enumerate}
Both families are orthonormal in $L^2[0,1]$.
\end{lemma}

\begin{proof}
Orthonormality is a direct computation with $2\cos a\cos b=\cos(a-b)+\cos(a+b)$. The inner products are $\langle\mathbf 1_{[0,s]},c_n\rangle=\sqrt2\frac{\sin n\pi s}{n\pi}$ and
$\langle\mathbf 1_{[0,s]},k_n\rangle=\sqrt2\frac{\sin((n-\frac12)\pi s)}{(n-\frac12)\pi}$, which gives the middle expressions.

(a) Let $f$, $g$ be the functions on $\R$ of period $2$, even, equal to $\mathbf 1_{[0,s]}$ and $\mathbf 1_{[0,t]}$ on $[0,1]$. They are piecewise continuous; after the change of variable $u=x/\pi$, Parseval reads
$\frac12\int_{-1}^1fg=\sum_kc_k(f)c_k(g)$ with $c_k(f)=\frac12\int_{-1}^1f(u)e^{-ik\pi u}\dd u=\int_0^1\mathbf 1_{[0,s]}(u)\cos k\pi u\dd u$ by evenness. The left side is $\min(s,t)$. On the right, $c_0(f)c_0(g)=st$ and the terms $k$ and $-k$
together give $2\langle\mathbf 1_{[0,s]},\cos n\pi\cdot\rangle\langle\mathbf 1_{[0,t]},\cos n\pi\cdot\rangle=\langle\mathbf 1_{[0,s]},c_n\rangle\langle\mathbf 1_{[0,t]},c_n\rangle$.

(b) is proved in the same way, with an extension of period $4$ that is odd about $u=1$, so that only the frequencies $(n-\frac12)\pi$ survive.
\end{proof}

\begin{proof}[Proof of (b)]
Let $f$ be the function on $\R$ of period $4$ with $f(u)=\mathbf 1_{[0,s]}(u)$ for $0\le u\le1$, $f(u)=-f(2-u)$ for $1<u\le2$, and $f(-u)=f(u)$. It is even and piecewise continuous, and its
coefficients are $c_k(f)=\frac12\int_0^2f(u)\cos\frac{k\pi u}2\dd u$. On $[1,2]$, the substitution $u=2-v$ gives $\int_1^2f(u)\cos\frac{k\pi u}2\dd u=-(-1)^k\int_0^1f(v)\cos\frac{k\pi v}2\dd v$, so
$c_k(f)=0$ for even $k$, and for odd $k=2n-1$,
\[c_{2n-1}(f)=\int_0^1\mathbf 1_{[0,s]}(u)\cos\bigl((n-\tfrac12)\pi u\bigr)\dd u=\tfrac1{\sqrt2}\langle\mathbf 1_{[0,s]},k_n\rangle .\]
With $g$ built from $t$ in the same way, $fg=\mathbf 1_{[0,s]}\mathbf 1_{[0,t]}$ on $[0,1]$ and on its three images in $[-2,2]$, so $\frac14\int_{-2}^2fg=\min(s,t)$. On the right of Parseval, the odd
indices $k=\pm(2n-1)$ contribute $2\cdot\frac12\langle\mathbf 1_{[0,s]},k_n\rangle\langle\mathbf 1_{[0,t]},k_n\rangle$.
\end{proof}

Uniform convergence of the series needs two tools. The first is an inequality of L\'evy for symmetric sums: the maximum of the partial sums is controlled by the last one. It holds for random elements of
$C[0,1]$ with the sup norm $\|\cdot\|$; a \emph{random continuous function} is a process with continuous paths, and then $\|X\|=\sup_{t\in\Q\cap[0,1]}|X_t|$ is a random variable.

\begin{lemma}[L\'evy's maximal inequality]\label{lem:levymax}
Let $Y_1,\dots,Y_N$ be independent random continuous functions on $[0,1]$, each symmetric ($-Y_j$ has the law of $Y_j$), and $S_k=Y_1+\dots+Y_k$. For every $x>0$,
\[
P\Bigl(\max_{k\le N}\|S_k\|>x\Bigr)\le2P(\|S_N\|>x).
\]
\end{lemma}

\begin{proof}
Let $\tau$ be the first $k$ with $\|S_k\|>x$, and $R_k=S_N-S_k$. Since $S_k=\frac12\bigl((S_k+R_k)+(S_k-R_k)\bigr)$, on $\{\tau=k\}$ either $\|S_k+R_k\|=\|S_N\|>x$ or $\|S_k-R_k\|>x$. The vector
$(Y_1,\dots,Y_k,-Y_{k+1},\dots,-Y_N)$ has the law of $(Y_1,\dots,Y_N)$, by independence and symmetry; the event $\{\tau=k\}$ depends only on $Y_1,\dots,Y_k$, and $S_k-R_k$ is the $S_N$ of the flipped vector. Hence
$P(\tau=k,\|S_k-R_k\|>x)=P(\tau=k,\|S_N\|>x)$, and
\[P(\tau=k)\le P(\tau=k,\|S_N\|>x)+P(\tau=k,\|S_k-R_k\|>x)=2P(\tau=k,\|S_N\|>x).\]
Sum over $k\le N$.
\end{proof}

The second tool bounds the maximum of many Gaussians, which grows only like the square root of the logarithm of their number.

\begin{lemma}[maxima of Gaussians]\label{lem:gaussmax}
If $G_1,\dots,G_M$ are centred Gaussians (not necessarily independent) with variances at most $\sigma^2$, then $\mathbb E\max_i|G_i|\le\sigma\sqrt{2\log(2M)}$.
\end{lemma}

\begin{proof}
For $\lambda>0$, $\mathbb E e^{\lambda|G_i|}\le\mathbb E e^{\lambda G_i}+\mathbb E e^{-\lambda G_i}\le2e^{\lambda^2\sigma^2/2}$. By Jensen's inequality for the convex function $\exp$,
\[
e^{\lambda\mathbb E\max_i|G_i|}\le\mathbb E e^{\lambda\max_i|G_i|}\le\sum_i\mathbb E e^{\lambda|G_i|}\le2Me^{\lambda^2\sigma^2/2},
\]
so $\mathbb E\max|G_i|\le\frac{\log(2M)}\lambda+\frac{\lambda\sigma^2}2$. Take $\lambda=\sqrt{2\log(2M)}/\sigma$.
\end{proof}

\begin{theorem}[Paley--Wiener--Zygmund; the Karhunen--Lo\`eve series]\label{thm:wienerseries}
Let $(\xi_n)_{n\ge0}$ be independent standard Gaussians. Almost surely, the series
\[
W_t=\xi_0t+\sqrt2\sum_{n\ge1}\xi_n\frac{\sin n\pi t}{n\pi}\qquad\text{and}\qquad V_t=\sqrt2\sum_{n\ge1}\xi_n\frac{\sin((n-\frac12)\pi t)}{(n-\frac12)\pi}
\]
converge uniformly on $[0,1]$, and their sums $(W_t)_{t\le1}$ and $(V_t)_{t\le1}$ are Brownian motions.
\end{theorem}

\begin{proof}
Both series have the form $\sum_n\xi_nf_n(t)$ with $f_n(t)=a_n\sin(b_nt)$, where $|a_n|\le\frac2{n}$ and $0<b_n\le\pi n$ for $n\ge1$ (the term $\xi_0t$ is harmless). Group the terms in dyadic blocks
\[
D_k(t)=\sum_{2^k\le n<2^{k+1}}\xi_nf_n(t),\qquad k\ge0.
\]
\emph{Size of a block.} Let $h=2^{-3k}$ and let $G$ be the grid $\{jh:0\le j\le2^{3k}\}$. Every $t\in[0,1]$ is within $h$ of a grid point, and $|D_k'(u)|\le\sum_{\text{block}}|\xi_n||a_n|b_n\le2\pi\sum_{\text{block}}|\xi_n|$, so
\[
\|D_k\|\le\max_{g\in G}|D_k(g)|+2\pi h\sum_{2^k\le n<2^{k+1}}|\xi_n|.
\]
Each $D_k(g)$ is centred Gaussian with variance $\sum_{\text{block}}a_n^2\sin^2(b_ng)\le\sum_{n\ge2^k}\frac4{n^2}\le8\cdot2^{-k}$, and $G$ has at most $2^{3k+1}$ points. By Lemma~\ref{lem:gaussmax}, and $\mathbb E|\xi_n|\le1$,
\[
\mathbb E\|D_k\|\le\sqrt{8\cdot2^{-k}}\sqrt{2\log(2^{3k+2})}+2\pi2^{-3k}2^k\le c\sqrt k\,2^{-k/2}
\]
for $k\ge1$ and a constant $c$.

\emph{Inside a block.} Each term $\xi_nf_n$ is a symmetric random continuous function, and the terms are independent. By Lemma~\ref{lem:levymax} and Markov's inequality,
\[
P\Bigl(\max_{2^k\le m<2^{k+1}}\Bigl\|\sum_{2^k\le n\le m}\xi_nf_n\Bigr\|>2^{-k/4}\Bigr)\le2P(\|D_k\|>2^{-k/4})\le2c\sqrt k\,2^{-k/4},
\]
which is summable in $k$. By Borel--Cantelli, almost surely, for all $k$ large, every partial sum inside the $k$-th block has norm at most $2^{-k/4}$. Then for $M<M'$ in blocks $k\le k'$, with $k$ large,
$\bigl\|\sum_{M<n\le M'}\xi_nf_n\bigr\|\le2\sum_{j\ge k}2^{-j/4}$, which tends to $0$: the partial sums form a Cauchy sequence in $C[0,1]$, and converge uniformly to a continuous function.

\emph{The limit is a Brownian motion.} For each $t$, the partial sums also converge in $L^2$, as noted above, so the uniform limit is their $L^2$ limit almost surely. The vectors of partial sums at $t_1,\dots,t_k$ are centred
Gaussian; their characteristic functions $e^{-u^{\mathsf T}\Sigma_Nu/2}$ converge to that of the limit (as $|e^{ia}-e^{ib}|\le|a-b|$ and $L^2$ convergence implies convergence in mean), and the covariance matrices
$\Sigma_N$ converge to the covariances of the limit, which are $\min(s,t)$ by Lemma~\ref{lem:twobases}. So the limit is a centred Gaussian process with covariance $\min(s,t)$ and continuous paths, a Brownian motion
on $[0,1]$ (Tome~3, Proposition~11.2).
\end{proof}

The second series is the Karhunen--Lo\`eve expansion of Tome~7, Chapter~4, whose uniform convergence was promised there. It is the expansion that converges fastest in mean square: no other orthonormal
family captures more of the variance with its first $N$ terms (Tome~7). Figure~\ref{fig:wienerseries} shows how few terms already give a convincing Brownian path.

\begin{figure}[ht]
\centering
\fig{ch01-wienerseries}
\caption{Partial sums of the Karhunen--Lo\`eve series of Theorem~\ref{thm:wienerseries}, with the same Gaussian coefficients: $N=4$ (smooth, dashed), $N=32$, and $N=1024$ terms. The first mode alone carries
$\frac8{\pi^2}\approx81\%$ of the variance; the rest is the roughness, spread over all frequencies.}
\label{fig:wienerseries}
\end{figure}

\begin{remark}
Kiyosi It\^o and Makiko Nisio proved in 1968 that the series $\sum_n\xi_n\int_0^te_n$ converges uniformly almost surely for \emph{every} complete orthonormal family $(e_n)$ of $L^2[0,1]$, in whatever order the terms are taken.
Their theorem is the infinite-dimensional version of the following fact, which the proof above already used in a special case: for sums of independent symmetric random vectors in a separable Banach space,
convergence in probability, almost sure convergence, and convergence in law are equivalent. L\'evy's own construction of Tome~3 is the same series for the family of \emph{Haar functions}, which makes uniform
convergence easy because the functions of one level have disjoint supports (Exercise~\ref{ex:haar}).
\end{remark}

\section{The Wiener integral}\label{sec:wienerint}

The series $W_t=\sum\xi_n\langle\mathbf 1_{[0,t]},e_n\rangle$ suggests a more general object: for $f\in L^2$, the Gaussian variable $\sum\xi_n\langle f,e_n\rangle$, which deserves the name $\int f\dd B$. We define it
directly from a Brownian motion. For a step function $f=\sum_{j=1}^ka_j\mathbf 1_{]t_{j-1},t_j]}$ with $0\le t_0<\dots<t_k$, let
\[
I(f)=\sum_{j=1}^ka_j(B_{t_j}-B_{t_{j-1}}).
\]
This does not depend on how $f$ is written (refining the subdivision does not change the sum), it is linear in $f$, and since the increments are independent with variances $t_j-t_{j-1}$, $I(f)$ is centred
Gaussian with $\mathbb E I(f)^2=\sum a_j^2(t_j-t_{j-1})=\|f\|_2^2$. By polarization $\mathbb E I(f)I(g)=\langle f,g\rangle$.

\begin{theorem}[the Wiener integral]\label{thm:wienerint}
The map $I$ extends uniquely to a linear isometry $f\mapsto\int_0^\infty f\dd B$ from $L^2[0,\infty[$ into $L^2(P)$. Every finite family $\bigl(\int f_1\dd B,\dots,\int f_k\dd B\bigr)$ is a centred Gaussian vector with
covariances $\langle f_j,f_l\rangle$. In particular $\int f\dd B$ and $\int g\dd B$ are independent when $f\perp g$.
\end{theorem}

\begin{proof}
Step functions are dense in $L^2[0,\infty[$ (reminder). An isometry defined on a dense subspace sends Cauchy sequences to Cauchy sequences, so it extends uniquely and continuously, by completeness of $L^2(P)$
(Tome~6, Theorem~9.4), and the extension is still a linear isometry. For step functions $f_1,\dots,f_k$ the vector $(I(f_j))_j$ is Gaussian, being a linear image of finitely many independent increments; for general
$f_j$, approximate each by step functions in $L^2$ and pass to the limit in the characteristic functions, as in the proof of Theorem~\ref{thm:wienerseries}. Uncorrelated components of a Gaussian vector are
independent.
\end{proof}

So $\int_0^tf\dd B$ is a Gaussian process with independent increments, of variance $\int_0^tf^2$, and one can check that it has a continuous modification (Exercise~\ref{ex:ou}). Two classical consequences:
integration by parts, $\int_0^1t\dd B_t=B_1-\int_0^1B_t\dd t$ (Exercise~\ref{ex:intbyparts}), and the \emph{Brownian motion on the whole line}: take two independent Brownian motions $B^+$, $B^-$ and set
$B_t=B^+_t$ for $t\ge0$, $B_t=B^-_{-t}$ for $t<0$. For $f\in L^2(\R)$, $\int_\R f\dd B=\int_0^\infty f(u)\dd B^+_u+\int_0^\infty f(-u)\dd B^-_u$ is again an isometry, from $L^2(\R)$, and its law is invariant under
translations of $f$: the covariances $\langle f(\cdot+a),g(\cdot+a)\rangle=\langle f,g\rangle$ do not change. This invariance, together with scaling, is what the next section exploits.

\begin{beauty}
The Wiener integral makes Brownian motion into a geometric object: the map $f\mapsto\int f\dd B$ is an isometric copy of the Hilbert space $L^2$ inside the space of random variables, and every Gaussian question
about $B$ becomes a question about angles and lengths of functions. Independence is orthogonality, the covariance is the inner product, and the three constructions of this chapter are three choices of orthonormal
basis for one and the same object.
\end{beauty}

\section{Fractional Brownian motion}\label{sec:fbm}

In 1951 the British hydrologist Harold Edwin Hurst, who had spent decades measuring the Nile to plan its dams, published a puzzling statistic. Over a period of length $n$, the range of the cumulated deviations of
the yearly river flow from its mean, divided by the standard deviation, grew like $n^H$ with $H$ near $0.73$, not like $n^{1/2}$ as independent years would give. Wet years cluster, and so do dry ones, over spans of
decades. In 1968 Beno\^it Mandelbrot and John Van Ness explained Hurst's law by a Gaussian process with stationary increments that scales like $t^H$. Kolmogorov had already met it in 1940, studying turbulence.

\begin{theorem}[fractional Brownian motion]\label{thm:fbm}
Let $0<H<1$. There is a centred Gaussian process $(X_t)_{t\ge0}$ with continuous paths and covariance
\[
\operatorname{Cov}(X_s,X_t)=\tfrac12\bigl(s^{2H}+t^{2H}-|t-s|^{2H}\bigr).
\]
It has stationary increments ($X_{t+h}-X_h$ has the law of $X_t$, as a process in $t$), it is self-similar ($c^{-H}X_{ct}$ has the law of $X_t$, as a process), and its paths are H\"older continuous of every order $\gamma<H$.
For $H=\frac12$ it is Brownian motion.
\end{theorem}

\begin{proof}
Let $B$ be a Brownian motion on the whole line, and for $t\ge0$ put
\[
g_t(u)=(t-u)_+^{H-1/2}-(-u)_+^{H-1/2},\qquad u\in\R,
\]
where $x_+=\max(x,0)$ and $0^{H-1/2}$ is read as $0$ (for $H=\frac12$, $g_t=\mathbf 1_{]0,t]}$). The function $g_t$ is in $L^2(\R)$: near $u=t$ and $u=0$ its square behaves like $|u-t|^{2H-1}$ or $|u|^{2H-1}$, integrable since
$2H-1>-1$; for $u\to-\infty$, the mean value theorem gives $|g_t(u)|\le|H-\frac12|\,t\,|u|^{H-3/2}$, and $|u|^{2H-3}$ is integrable at infinity since $2H-3<-1$. Let $Y_t=\int_\R g_t\dd B$, a centred Gaussian process
(Theorem~\ref{thm:wienerint}).

\emph{Stationary increments.} $g_{t+h}-g_h=g_t(\cdot-h)$, a translate. Since the covariances of a Wiener integral are inner products, which translations preserve, the process $(Y_{t+h}-Y_h)_t$ has the same covariances as
$(Y_t)_t$, hence the same law. \emph{Scaling.} $g_{ct}(u)=c^{H-1/2}g_t(u/c)$, so $\|g_{ct}\|^2=c^{2H-1}\cdot c\|g_t\|^2=c^{2H}\|g_t\|^2$, and more generally $\langle g_{cs},g_{ct}\rangle=c^{2H}\langle g_s,g_t\rangle$. Hence
$\operatorname{Var}Y_t=\kappa^2t^{2H}$ with $\kappa=\|g_1\|>0$, and $\operatorname{Var}(Y_t-Y_s)=\operatorname{Var}Y_{|t-s|}=\kappa^2|t-s|^{2H}$. Polarization,
$\operatorname{Cov}(Y_s,Y_t)=\frac12\bigl(\operatorname{Var}Y_s+\operatorname{Var}Y_t-\operatorname{Var}(Y_t-Y_s)\bigr)$, gives $\kappa^2$ times the announced covariance. Put $X=Y/\kappa$.

\emph{Continuity.} $\mathbb E|X_t-X_s|^{2k}=(2k-1)!!\,|t-s|^{2Hk}$, so Theorem~\ref{thm:kolmcont} with $p=2k$, $\beta=2Hk-1$ (positive for $k$ large) gives a continuous modification, H\"older of every order
$<H-\frac1{2k}$, on each $[j,j+1]$; let $k\to\infty$ and glue as in Corollary~\ref{cor:bmholder}. The modification does not change the finite-dimensional laws, so stationarity and scaling hold for it. For $H=\frac12$ the
covariance is $\frac12(s+t-|t-s|)=\min(s,t)$.
\end{proof}

\begin{think}
For $H\neq\frac12$, are the increments over disjoint intervals positively or negatively correlated? Compute $\operatorname{Cov}(X_1-X_0,X_2-X_1)$ before reading on.
\end{think}

Polarization gives $\operatorname{Cov}(X_1,X_2-X_1)=\frac12(2^{2H}-2)$, positive for $H>\frac12$ and negative for $H<\frac12$. More generally the correlation of $X_1-X_0$ with $X_{n+1}-X_n$ behaves like
$H(2H-1)n^{2H-2}$ (Exercise~\ref{ex:fbmcorr}). For $H>\frac12$ the correlations decay so slowly that they are not summable: the process has \emph{long memory}, Hurst's Nile. For $H<\frac12$ it is anti-persistent, each
step tending to undo the last, and its paths are rougher than Brownian ones (Figure~\ref{fig:fbm}).

\begin{figure}[ht]
\centering
\fig{ch01-fbm}
\caption{Fractional Brownian motions on $[0,1]$ with $H=0.25$, $0.5$ (Brownian motion) and $0.75$, sampled exactly at $512$ points from their covariance (Cholesky factorization). Rough and anti-persistent, then
Brownian, then smooth and persistent.}
\label{fig:fbm}
\end{figure}

\begin{remark}
Fractional Brownian motion is not a martingale and, for $H\neq\frac12$, its quadratic variation along fine partitions is $0$ or infinite (Exercise~\ref{ex:fbmqv}). The stochastic calculus of Chapters~3--4, built on
martingales, does not apply to it, which is one reason why Brownian motion, with $H=\frac12$ exactly, is the canonical random motion.
\end{remark}

\section{Exercises}

\begin{exercise}[$\star$]\label{ex:notevent}
Let $X_t=0$ and $Y_t=\mathbf 1_{\{t=U\}}$ for $t\in[0,1]$, with $U$ uniform on $[0,1]$. Show that $Y$ is a modification of $X$, that $P(Y\text{ is continuous})=0$, and that the set $\{x:\sup_tx(t)\ge1\}$ is not in the product $\sigma$-algebra of $\R^{[0,1]}$. So the law of a
process on $\R^{[0,1]}$ cannot tell whether its paths reach $1$.
\end{exercise}

\begin{exercise}[$\star$]\label{ex:bridgeseries}
Show that $W_t-tW_1=\sqrt2\sum_{n\ge1}\xi_n\frac{\sin n\pi t}{n\pi}$ for the series $W$ of Theorem~\ref{thm:wienerseries}, and deduce the expansion of the Brownian bridge (Example~\ref{ex:kernels}(b)) as a random
Fourier sine series. Use it to show that $\int_0^1(B_t-tB_1)^2\dd t$ has the law of $\sum_{n\ge1}\frac{\xi_n^2}{n^2\pi^2}$, and compute its mean.
\end{exercise}

\begin{exercise}[$\star$]\label{ex:intbyparts}
Show that $\int_0^1t\dd B_t=B_1-\int_0^1B_t\dd t$, where the last integral is a Riemann integral of the continuous path. Deduce that $\int_0^1B_t\dd t$ is $N(0,\frac13)$, and compute its covariance with $B_1$.
\end{exercise}

\begin{exercise}[$\star\star$]\label{ex:sheet}
\emph{The Brownian sheet.} Show that $K((s_1,s_2),(t_1,t_2))=\min(s_1,t_1)\min(s_2,t_2)$ is positive semi-definite on $[0,1]^2$, and that the Gaussian process with this covariance has a continuous modification,
H\"older of every order $<\frac12$. (Bound $\mathbb E(X_t-X_s)^2$ by $2|t-s|$ in the max norm.) Show that along each horizontal line $s_2=c>0$, $X_{(\cdot,c)}/\sqrt c$ is a Brownian motion.
\end{exercise}

\begin{exercise}[$\star\star$]\label{ex:ou}
Let $f\in L^2_{\mathrm{loc}}[0,\infty[$ (square integrable on bounded intervals). Show that $M_t=\int_0^tf\dd B$ has a continuous modification. (Write $M_t=W_{A(t)}$ in law with $A(t)=\int_0^tf^2$, or apply
Theorem~\ref{thm:kolmcont} after a time change.) Then show that $X_t=e^{-t}\bigl(X_0+\sqrt2\int_0^te^u\dd B_u\bigr)$, with $X_0\sim N(0,1)$ independent of $B$, is a centred Gaussian process with covariance $e^{-|t-s|}$:
the stationary Ornstein--Uhlenbeck process of Example~\ref{ex:kernels}(c).
\end{exercise}

\begin{exercise}[$\star\star$]\label{ex:haar}
\emph{L\'evy's construction as a series.} The Haar functions on $[0,1]$ are $h_0=1$ and $h_{n,k}=2^{n/2}\bigl(\mathbf 1_{[2k2^{-n-1},(2k+1)2^{-n-1}[}-\mathbf 1_{[(2k+1)2^{-n-1},(2k+2)2^{-n-1}[}\bigr)$ for $n\ge0$, $0\le k<2^n$.
Their integrals $s_{n,k}(t)=\int_0^th_{n,k}$ are tents of height $2^{-n/2-1}$ with disjoint supports for fixed $n$. Show that $\sum_{n,k}\xi_{n,k}s_{n,k}$ converges uniformly almost surely, using only
$\|\sum_k\xi_{n,k}s_{n,k}\|=2^{-n/2-1}\max_k|\xi_{n,k}|$. (That the limit is a Brownian motion would follow from completeness of the Haar system, which you may admit; Tome~3 proved it differently.)
\end{exercise}

\begin{exercise}[$\star\star$]\label{ex:fbmcorr}
For fractional Brownian motion with index $H$, show that $\rho(n)=\operatorname{Cov}(X_1-X_0,X_{n+1}-X_n)=\frac12\bigl((n+1)^{2H}-2n^{2H}+(n-1)^{2H}\bigr)$ for $n\ge1$, and that $\rho(n)\sim H(2H-1)n^{2H-2}$ as $n\to\infty$. For
which $H$ is $\sum_n|\rho(n)|$ finite? Show that $\operatorname{Var}\bigl(\sum_{j=1}^n(X_j-X_{j-1})\bigr)=n^{2H}$, and explain Hurst's observation.
\end{exercise}

\begin{exercise}[$\star\star\star$]\label{ex:fbmqv}
Let $X$ be a fractional Brownian motion of index $H$ and $Q_n=\sum_{j=1}^n(X_{j/n}-X_{(j-1)/n})^2$. Show that $\mathbb E Q_n=n^{1-2H}$. For $H>\frac12$ deduce that $Q_n\to0$ in $L^1$. For $H<\frac12$, show that
$\operatorname{Var}Q_n=2n^{-4H}\sum_{j,k\le n}\rho(|j-k|)^2$ with $\rho(0)=1$ and $\rho$ as in Exercise~\ref{ex:fbmcorr}, that this is $O(n^{1-4H})$, and deduce that $Q_n\to\infty$ in probability.
\end{exercise}

\section{Hints}

\begin{hint}{ex:notevent}
For the last part, every event of the product $\sigma$-algebra is determined by countably many coordinates (proof of Theorem~\ref{thm:kolmext}).
\end{hint}

\begin{hint}{ex:bridgeseries}
$\xi_0t$ is the only term not vanishing at $t=1$, and every term $\sin n\pi t$ vanishes there. For the integral, use the orthogonality of the sines on $[0,1]$.
\end{hint}

\begin{hint}{ex:intbyparts}
Check it for the step function approximations $\sum_j t_{j-1}\mathbf 1_{]t_{j-1},t_j]}$ of $t$, by summation by parts, and pass to the limit.
\end{hint}

\begin{hint}{ex:sheet}
The kernel is $\int\mathbf 1_{[0,s_1]\times[0,s_2]}\mathbf 1_{[0,t_1]\times[0,t_2]}$, an inner product in $L^2([0,1]^2)$.
\end{hint}

\begin{hint}{ex:ou}
For the covariance of $X$, compute $\mathbb E X_sX_t$ with Theorem~\ref{thm:wienerint}, for $s\le t$.
\end{hint}

\begin{hint}{ex:haar}
Use $P(\max_{k<2^n}|\xi_{n,k}|>2\sqrt n)\le2^ne^{-2n}$.
\end{hint}

\begin{hint}{ex:fbmcorr}
Write the increments' covariance with polarization: $2\operatorname{Cov}(a,b)=\operatorname{Var}(a+b)-\operatorname{Var}a-\operatorname{Var}b$, with $a=X_1-X_0$ and $b=X_{n+1}-X_n$ grouped as differences of $X$ over
intervals of lengths $n+1$, $n$ and $n-1$. For the asymptotics, it is a second difference of $x^{2H}$.
\end{hint}

\begin{hint}{ex:fbmqv}
For centred jointly Gaussian $U,V$, $\operatorname{Cov}(U^2,V^2)=2\operatorname{Cov}(U,V)^2$. Then use $\sum_{k<n}\rho(k)^2=O(1)$ when $4H-4<-1$.
\end{hint}

\section{Solutions}

\begin{solution}{ex:notevent}
\idea The two processes differ only at one random time, which no fixed time can catch; but a supremum over all times does catch it, and such a supremum looks at uncountably many coordinates.
\writeup For fixed $t$, $P(Y_t\neq X_t)=P(U=t)=0$, so $Y$ is a modification of $X$. Every path of $Y$ jumps from $0$ to $1$ at $t=U$, so $P(Y\text{ continuous})=0$. Suppose
$E=\{x:\sup_tx(t)\ge1\}$ were in the product $\sigma$-algebra; then $E=\pi_S^{-1}(A)$ for a countable $S$ and some $A$. Pick $t_0\in[0,1]\setminus S$. The zero function and the indicator of $\{t_0\}$ agree on $S$, so
either both are in $E$ or neither is; but the indicator is in $E$ and $0$ is not. Since $X$ and $Y$ have the same law on $\R^{[0,1]}$ while $\sup X=0$ and $\sup Y=1$, the law alone indeed cannot see the supremum: continuity and
suprema must come from a good choice of modification.
\end{solution}

\begin{solution}{ex:bridgeseries}
\idea The bridge is what remains after removing the straight line from $0$ to $B_1$; in the series, the line is exactly the first term.
\writeup At $t=1$ all the sines vanish, so $W_1=\xi_0$ and $W_t-tW_1=\sqrt2\sum_{n\ge1}\xi_n\frac{\sin n\pi t}{n\pi}$, uniformly convergent. Since $W$ is a Brownian motion, the left side is a Brownian bridge. The functions
$\sqrt2\sin n\pi t$ are orthonormal on $[0,1]$, and the series converges uniformly, so we may integrate the square term by term (the cross terms vanish by orthogonality):
$\int_0^1(W_t-tW_1)^2\dd t=\sum_{n\ge1}\frac{\xi_n^2}{n^2\pi^2}$. Its mean is $\sum\frac1{n^2\pi^2}=\frac16$, which agrees with $\int_0^1t(1-t)\dd t=\frac16$.
\end{solution}

\begin{solution}{ex:intbyparts}
\idea Summation by parts is integration by parts for step functions; both sides are continuous in $L^2$, so the identity survives the limit.
\writeup Let $0=t_0<\dots<t_m=1$ with mesh $\to0$, and $f_m=\sum_jt_{j-1}\mathbf 1_{]t_{j-1},t_j]}$, which tends to $t\mapsto t$ in $L^2[0,1]$. Summation by parts:
$\sum_jt_{j-1}(B_{t_j}-B_{t_{j-1}})=B_1-\sum_jB_{t_j}(t_j-t_{j-1})$. As the mesh tends to $0$, the left side tends to $\int_0^1t\dd B_t$ in $L^2$ by the isometry, and the Riemann sum on the right tends to
$\int_0^1B_t\dd t$ surely, by continuity of the path; limits in $L^2$ and almost sure limits agree almost surely. So $\int_0^1B_t\dd t=\int_0^1(1-t)\dd B_t$, a Wiener integral, hence Gaussian, centred, with variance
$\int_0^1(1-t)^2\dd t=\frac13$, and its covariance with $B_1=\int_0^11\dd B$ is $\int_0^1(1-t)\dd t=\frac12$.
\end{solution}

\begin{solution}{ex:sheet}
\idea Positive semi-definiteness is easiest when the kernel is an inner product of indicators; continuity comes from Kolmogorov in dimension $d=2$ with high moments.
\writeup With $R_s=[0,s_1]\times[0,s_2]$, $K(s,t)=\lambda(R_s\cap R_t)=\langle\mathbf 1_{R_s},\mathbf 1_{R_t}\rangle$, so $\sum c_jc_kK(t_j,t_k)=\|\sum c_j\mathbf 1_{R_{t_j}}\|^2\ge0$, and the process exists by
Corollary~\ref{cor:gaussexist}. Then $\mathbb E(X_t-X_s)^2=\|\mathbf 1_{R_t}-\mathbf 1_{R_s}\|^2=\lambda(R_t\triangle R_s)$, and the symmetric difference of two rectangles in $[0,1]^2$ with corners at max-distance $\delta$ has
area at most $2\delta$. Gaussianity gives $\mathbb E|X_t-X_s|^{2k}\le(2k-1)!!(2|t-s|)^k$, so Theorem~\ref{thm:kolmcont} with $d=2$, $p=2k$, $\beta=k-2$ gives H\"older continuity of every order $<\frac{k-2}{2k}$, hence of
every order $<\frac12$. On the line $s_2=c$, $\operatorname{Cov}(X_{(s,c)},X_{(t,c)})=c\min(s,t)$, so $X_{(\cdot,c)}/\sqrt c$ is a centred Gaussian process with continuous paths and covariance $\min(s,t)$: a Brownian motion.
\end{solution}

\begin{solution}{ex:ou}
\idea A Wiener integral $\int_0^tf\dd B$ has independent increments of variance $\int_s^tf^2$: it is a Brownian motion run with the clock $A(t)$, and Kolmogorov's criterion sees this through its moments.
\writeup Fix $T$ and let $A(t)=\int_0^tf^2$, continuous and non-decreasing. Then $\mathbb E|M_t-M_s|^{2k}=(2k-1)!!(A(t)-A(s))^k$. Run the process with the clock $A$: for $a\in[0,A(T)]$ let
$\tau(a)=\min\{t:A(t)\ge a\}$, so that $A(\tau(a))=a$, and $Z_a=M_{\tau(a)}$. Then $\mathbb E|Z_a-Z_b|^{2k}=(2k-1)!!|a-b|^k$, and Theorem~\ref{thm:kolmcont} gives a continuous modification $\tilde Z$. Put
$\tilde M_t=\tilde Z_{A(t)}$, continuous in $t$. For fixed $t$, $\tilde M_t=Z_{A(t)}=M_{\tau(A(t))}$ almost surely, and $M_{\tau(A(t))}-M_t$ has variance $A(t)-A(\tau(A(t)))=0$; so $\tilde M$ is a modification of $M$.
For the Ornstein--Uhlenbeck process, with $s\le t$ and the independence of $X_0$ and $B$,
\[
\mathbb E X_sX_t=e^{-s-t}\Bigl(1+2\int_0^se^{2u}\dd u\Bigr)=e^{-s-t}\bigl(1+e^{2s}-1\bigr)=e^{-(t-s)}.
\]
\end{solution}

\begin{solution}{ex:haar}
\idea At level $n$ the tents do not overlap, so the norm of the level is the largest coefficient times the height of one tent; the largest of $2^n$ Gaussians is about $\sqrt{2n\log2}$, tiny against $2^{n/2}$.
\writeup For fixed $n$, the supports of $s_{n,k}$ are disjoint, so $\|\sum_k\xi_{n,k}s_{n,k}\|=2^{-n/2-1}\max_k|\xi_{n,k}|$. By a union bound and $P(|\xi|>x)\le e^{-x^2/2}$ for $x\ge1$,
$P(\max_k|\xi_{n,k}|>2\sqrt n)\le2^ne^{-2n}$, summable. By Borel--Cantelli, almost surely $\max_k|\xi_{n,k}|\le2\sqrt n$ for $n$ large, and then $\sum_n\|\sum_k\xi_{n,k}s_{n,k}\|\le\sum_n\sqrt n2^{-n/2}<\infty$. The series of levels
converges absolutely in $C[0,1]$, uniformly.
\end{solution}

\begin{solution}{ex:fbmcorr}
\idea Everything about fractional Brownian motion follows from one formula: the variance of an increment over an interval of length $\ell$ is $\ell^{2H}$. Express any covariance of increments through variances of
increments, and the answer is a second difference.
\writeup By stationarity, $\operatorname{Var}(X_b-X_a)=|b-a|^{2H}$. With $a=X_1-X_0$, $b=X_{n+1}-X_n$: $a+b+(X_n-X_1)=X_{n+1}-X_0$, and the identity
$2\operatorname{Cov}(X_1-X_0,X_{n+1}-X_n)=\operatorname{Var}(X_{n+1}-X_0)+\operatorname{Var}(X_n-X_1)-\operatorname{Var}(X_{n}-X_0)-\operatorname{Var}(X_{n+1}-X_1)$ (expand all four in terms of $X_0,X_1,X_n,X_{n+1}$)
gives $\rho(n)=\frac12((n+1)^{2H}+(n-1)^{2H}-2n^{2H})$. With $\phi(x)=x^{2H}$, this is $\frac12(\phi(n+1)-2\phi(n)+\phi(n-1))=\frac12\phi''(\theta)$ for some $\theta\in(n-1,n+1)$, and
$\frac12\phi''(\theta)=H(2H-1)\theta^{2H-2}\sim H(2H-1)n^{2H-2}$. So $\sum|\rho(n)|<\infty$ exactly when $2H-2<-1$, i.e.\ $H<\frac12$ (and $\rho=0$ for $H=\frac12$). Finally the sum of the first $n$ increments is $X_n-X_0$,
of variance $n^{2H}$: the cumulated deviations of $n$ years grow like $n^H$, and with $H\approx0.73$ the years are positively correlated over long spans, which is what Hurst saw in the Nile.
\end{solution}

\begin{solution}{ex:fbmqv}
\idea The mean is immediate from scaling. For $H<\frac12$ the mean blows up, and one needs the variance to be small compared with the square of the mean, which the decay of correlations provides.
\writeup Each increment over $[\frac{j-1}n,\frac jn]$ has variance $n^{-2H}$, so $\mathbb E Q_n=n\cdot n^{-2H}$, which tends to $0$ when $H>\frac12$: $Q_n\to0$ in $L^1$. For centred jointly Gaussian $U,V$ with
$\operatorname{Cov}(U,V)=c$, write $V=\frac c{\operatorname{Var}U}U+V'$ with $V'$ independent of $U$; then $\operatorname{Cov}(U^2,V^2)=2c^2$. The normalized increments $n^H(X_{j/n}-X_{(j-1)/n})$ have, by scaling, the joint
law of the integer increments, whose covariances are $\rho(|j-k|)$. Hence $\operatorname{Var}Q_n=n^{-4H}\sum_{j,k\le n}2\rho(|j-k|)^2\le2n^{1-4H}\sum_{m\ge0}\rho(m)^2\cdot2$, and $\sum\rho(m)^2<\infty$ since
$\rho(m)^2=O(m^{4H-4})$ and $4H-4<-1$. Therefore $\operatorname{Var}Q_n/(\mathbb E Q_n)^2=O(n^{1-4H}/n^{2-4H})=O(1/n)\to0$, and by Chebyshev, $P(Q_n<\frac12\mathbb E Q_n)\to0$, while $\mathbb E Q_n=n^{1-2H}\to\infty$. So
$Q_n\to\infty$ in probability.
\end{solution}

\section*{Sources and further reading}

\begin{itemize}
\item \emph{History.} R.~Brown, ``A brief account of microscopical observations made in the months of June, July and August, 1827, on the particles contained in the pollen of plants; and on the general existence of
active molecules in organic and inorganic bodies'', \emph{Philosophical Magazine} 4 (1828), 161--173. A.~Einstein, ``\"Uber die von der molekularkinetischen Theorie der W\"arme geforderte Bewegung von in ruhenden
Fl\"ussigkeiten suspendierten Teilchen'', \emph{Annalen der Physik} 17 (1905), 549--560. J.~Perrin, ``Mouvement brownien et r\'ealit\'e mol\'eculaire'', \emph{Annales de Chimie et de Physique} (8) 18 (1909), 5--114
(English translation by F.~Soddy, \emph{Brownian Movement and Molecular Reality}, Taylor and Francis, 1910), and \emph{Les Atomes}, Alcan, Paris, 1913, whose preface contains the remark on curves without tangents
paraphrased in the opening. L.~Bachelier, ``Th\'eorie de la sp\'eculation'', \emph{Annales scientifiques de l'\'Ecole Normale Sup\'erieure} (3) 17 (1900), 21--86. N.~Wiener, ``Differential-space'', \emph{Journal of
Mathematics and Physics} 2 (1923), 131--174. R.~E.~A.~C.~Paley, N.~Wiener and A.~Zygmund, ``Notes on random functions'', \emph{Mathematische Zeitschrift} 37 (1933), 647--668. A.~N.~Kolmogorov, \emph{Grundbegriffe der
Wahrscheinlichkeitsrechnung}, Springer, 1933. E.~Slutsky, ``Alcune proposizioni sulla teoria delle funzioni aleatorie'', \emph{Giornale dell'Istituto Italiano degli Attuari} 8 (1937), 183--199 (the continuity criterion,
credited there to Kolmogorov). N.~N.~Chentsov, ``Weak convergence of stochastic processes whose trajectories have no discontinuities of the second kind'', \emph{Theory of Probability and its Applications} 1 (1956),
140--144. K.~Karhunen, ``\"Uber lineare Methoden in der Wahrscheinlichkeitsrechnung'', \emph{Annales Academiae Scientiarum Fennicae} A I 37 (1947). K.~It\^o and M.~Nisio, ``On the convergence of sums of independent Banach
space valued random variables'', \emph{Osaka Journal of Mathematics} 5 (1968), 35--48. H.~E.~Hurst, ``Long-term storage capacity of reservoirs'', \emph{Transactions of the American Society of Civil Engineers} 116 (1951),
770--799. B.~B.~Mandelbrot and J.~W.~Van Ness, ``Fractional Brownian motions, fractional noises and applications'', \emph{SIAM Review} 10 (1968), 422--437. A.~N.~Kolmogorov, ``Wienersche Spiralen und einige andere
interessante Kurven im Hilbertschen Raum'', \emph{Doklady Akad.\ Nauk SSSR} 26 (1940), 115--118.
\item J.-F.~Le Gall, \emph{Brownian Motion, Martingales, and Stochastic Calculus}, Springer, 2016, Chapters~1--2. Gaussian spaces, the Wiener integral and Kolmogorov's criterion, the route followed here.
\item P.~M\"orters and Y.~Peres, \emph{Brownian Motion}, Cambridge University Press, 2010, Chapter~1. Constructions and first properties, with the Paley--Wiener series as an exercise.
\item J.-P.~Kahane, \emph{Some Random Series of Functions}, 2nd ed., Cambridge University Press, 1985. Random Fourier series in depth, from L\'evy's inequality to the Salem--Zygmund theorems.
\item D.~Nualart, \emph{The Malliavin Calculus and Related Topics}, 2nd ed., Springer, 2006, Chapter~5. Fractional Brownian motion and its stochastic calculus.
\item S.~G.~Brush, ``A history of random processes I: Brownian movement from Brown to Perrin'', \emph{Archive for History of Exact Sciences} 5 (1968), 1--36. The physics story told with care.
\end{itemize}
